\documentclass[12pt,a4paper]{article}

\usepackage{fontawesome5}
\providecommand{\faPencilAlt}{\faIcon{pencil-alt}}

\usepackage{styles/couleurs}
\usepackage{styles/mise_en_page}
\usepackage{styles/boites}
\usepackage{styles/maths}
\usepackage{styles/figures}
\usepackage{styles/evaluations}

\usepackage{enumitem}
\usepackage{array}
\usepackage{tabularx}
\usepackage{amsmath}
\usepackage{amssymb}
\usepackage{tikz}
\usepackage{xparse}

\usetikzlibrary{calc,patterns}

\setlist[itemize]{
  leftmargin=1.4em,
  itemsep=0.15em,
  topsep=0.2em
}

\setlist[enumerate]{
  leftmargin=*,
  itemsep=0.25em,
  topsep=0.2em
}

% =========================================================
% EN-TÊTE
% =========================================================

\NewDocumentCommand{\entetestats}{}{%
\begin{tcolorbox}[
  enhanced,
  colback=NeutreBg,
  colframe=NeutreBg,
  boxrule=0pt,
  arc=2mm,
  borderline west={1.2mm}{0pt}{TitreC},
  left=4mm,
  right=4mm,
  top=3mm,
  bottom=3mm,
  before skip=0pt,
  after skip=4mm
]

\begin{center}
{\sffamily\bfseries\Large INTERROGATION -- 4e}

\vspace{0.15em}

{\sffamily\color{NeutreC}
Statistiques : vocabulaire, moyenne et médiane}
\end{center}

\vspace{0.65em}

\noindent
\begin{tabularx}{\linewidth}{
  @{}
  >{\raggedright\arraybackslash}X
  >{\raggedright\arraybackslash}X
  @{}
}
\textbf{Nom :}\ \rule{5.2cm}{0.4pt}
&
\textbf{Prénom :}\ \rule{5.2cm}{0.4pt}
\end{tabularx}

\vspace{0.7em}

\noindent
\begin{tabularx}{\linewidth}{
  @{}
  >{\raggedright\arraybackslash}X
  >{\raggedleft\arraybackslash}X
  @{}
}
{\small
\faCalculator\hspace{0.4em}
Calculatrice autorisée}
&
{\bfseries Note : \rule{1.5cm}{0.4pt}\,/\,10}
\end{tabularx}

\end{tcolorbox}
}

% =========================================================
% REMISE À ZÉRO DU NUMÉRO DES EXERCICES
% utile avant la version adaptée
% =========================================================

\makeatletter
\newcommand{\resetexos}{%
  \@ifundefined{c@exo}{}{\setcounter{exo}{0}}%
  \@ifundefined{c@exercice}{}{\setcounter{exercice}{0}}%
  \@ifundefined{c@exoeval}{}{\setcounter{exoeval}{0}}%
  \@ifundefined{c@exerciceeval}{}{\setcounter{exerciceeval}{0}}%
}
\makeatother


\begin{document}

% =========================================================
% =========================================================
% VERSION NORMALE
% =========================================================
% =========================================================

\entetestats

% =========================================================
% EXERCICE 1
% =========================================================

\exo{4}[Communiquer]

On a relevé le \textbf{temps mis pour venir au collège}
par les élèves d'une classe de 4e.

\begin{center}
\renewcommand{\arraystretch}{1.5}
\begin{tabular}{|c|c|c|c|c|c|}
\hline
Temps (en minutes)
& 5 & 10 & 15 & 20 & 25 \\
\hline
Effectif
& 4 & 9 & 6 & 3 & 3 \\
\hline
\end{tabular}
\end{center}

\begin{enumerate}

  \item Quelle est la \textbf{population étudiée} ? \pts{0,5}

  \lignes{1}

  \item Quel est le \textbf{caractère étudié} ? \pts{0,5}

  \lignes{1}

  \item Donne une \textbf{valeur} prise par ce caractère. \pts{0,5}

  \lignes{1}

  \item Quel est l'\textbf{effectif} correspondant à la valeur
  \(15\) minutes ? \pts{0,5}

  \lignes{1}

  \item Quel est l'\textbf{effectif total} de cette série ? \pts{1}

  \lignes{2}

  \item Combien d'élèves mettent \textbf{au moins 20 minutes}
  pour venir au collège ? \pts{1}

  \lignes{2}

\end{enumerate}

% =========================================================
% PAGE 2
% =========================================================

\newpage

% =========================================================
% EXERCICE 2
% =========================================================

\exo{3}[Calculer]

Voici les notes obtenues par sept élèves :

\[
9\qquad
12\qquad
14\qquad
11\qquad
16\qquad
13\qquad
9
\]

Calcule la \textbf{moyenne} de cette série.

Le calcul doit apparaître. \pts{3}

\lignes{6}

\vspace{1em}

% =========================================================
% EXERCICE 3
% =========================================================

\exo{3}[Calculer]

Détermine la \textbf{médiane} de la série suivante. \pts{3}

\[
14\qquad
7\qquad
11\qquad
9\qquad
16\qquad
12\qquad
8\qquad
13
\]

\lignes{7}


% =========================================================
% =========================================================
% VERSION ADAPTÉE
% =========================================================
% =========================================================

\newpage

\resetexos

\entetestats

% =========================================================
% EXERCICE 1
% =========================================================

\exo{3}[Communiquer]

On a relevé le \textbf{temps mis pour venir au collège}
par les élèves d'une classe de 4e.

\begin{center}
\renewcommand{\arraystretch}{1.5}
\begin{tabular}{|c|c|c|c|c|c|}
\hline
Temps (en minutes)
& 5 & 10 & 15 & 20 & 25 \\
\hline
Effectif
& 4 & 9 & 6 & 3 & 3 \\
\hline
\end{tabular}
\end{center}

\begin{enumerate}

  \item Quelle est la \textbf{population étudiée} ? \pts{0,5}

  \lignes{1}

  \item Quel est le \textbf{caractère étudié} ? \pts{0,5}

  \lignes{1}

  \item Donne une \textbf{valeur} prise par ce caractère. \pts{0,5}

  \lignes{1}

  \item Quel est l'\textbf{effectif} correspondant à la valeur
  \(15\) minutes ? \pts{0,5}

  \lignes{1}

  \item Quel est l'\textbf{effectif total} de cette série ? \pts{1}

  \lignes{2}

\end{enumerate}

% =========================================================
% PAGE 4
% =========================================================

\newpage

% =========================================================
% EXERCICE 2
% =========================================================

\exo{3}[Calculer]

Voici les notes obtenues par cinq élèves :

\[
9\qquad
12\qquad
14\qquad
11\qquad
14
\]

Calcule la \textbf{moyenne} de cette série.

Le calcul doit apparaître. \pts{3}

\lignes{7}

\vspace{1em}

% =========================================================
% EXERCICE 3
% =========================================================

\exo{4}[Calculer]

Détermine la \textbf{médiane} de la série suivante. \pts{4}

\[
14\qquad
8\qquad
11\qquad
9\qquad
13\qquad
10
\]

\lignes{8}


% =========================================================
% BARÈME ENSEIGNANT
% =========================================================
%
% =========================================================
% VERSION NORMALE
% =========================================================
%
% EXERCICE 1 : 4 points
%
% 1. Population étudiée : 0,5 pt
%    -> les élèves de la classe de 4e
%
% 2. Caractère étudié : 0,5 pt
%    -> le temps mis pour venir au collège
%
% 3. Une valeur du caractère : 0,5 pt
%    -> 5 ; 10 ; 15 ; 20 ou 25 minutes
%
% 4. Effectif correspondant à 15 min : 0,5 pt
%    -> 6 élèves
%
% 5. Effectif total : 1 pt
%    -> calcul 4+9+6+3+3 : 0,5 pt
%    -> réponse 25 élèves : 0,5 pt
%
% 6. Au moins 20 minutes : 1 pt
%    -> calcul 3+3 : 0,5 pt
%    -> réponse 6 élèves : 0,5 pt
%
%
% EXERCICE 2 : 3 points
%
% 9+12+14+11+16+13+9 = 84
% 84 / 7 = 12
%
% -> somme correcte : 1 pt
% -> division par 7 : 1 pt
% -> moyenne = 12 : 1 pt
%
%
% EXERCICE 3 : 3 points
%
% Série rangée :
% 7 ; 8 ; 9 ; 11 ; 12 ; 13 ; 14 ; 16
%
% Valeurs centrales : 11 et 12
%
% (11+12)/2 = 11,5
%
% -> rangement correct : 1 pt
% -> calcul avec les deux valeurs centrales : 1 pt
% -> médiane = 11,5 : 1 pt
%
% TOTAL : 10 points
%
%
% =========================================================
% VERSION ADAPTÉE
% =========================================================
%
% EXERCICE 1 : 3 points
%
% 1. Population étudiée : 0,5 pt
%    -> les élèves de la classe de 4e
%
% 2. Caractère étudié : 0,5 pt
%    -> le temps mis pour venir au collège
%
% 3. Une valeur du caractère : 0,5 pt
%    -> 5 ; 10 ; 15 ; 20 ou 25 minutes
%
% 4. Effectif correspondant à 15 min : 0,5 pt
%    -> 6 élèves
%
% 5. Effectif total : 1 pt
%    -> calcul 4+9+6+3+3 : 0,5 pt
%    -> réponse 25 élèves : 0,5 pt
%
%
% EXERCICE 2 : 3 points
%
% 9+12+14+11+14 = 60
% 60 / 5 = 12
%
% -> somme correcte : 1 pt
% -> division par 5 : 1 pt
% -> moyenne = 12 : 1 pt
%
%
% EXERCICE 3 : 4 points
%
% Série rangée :
% 8 ; 9 ; 10 ; 11 ; 13 ; 14
%
% Valeurs centrales : 10 et 11
%
% (10+11)/2 = 10,5
%
% -> rangement correct : 1,5 pt
% -> repérage des deux valeurs centrales : 1 pt
% -> calcul : 0,5 pt
% -> médiane = 10,5 : 1 pt
%
% TOTAL : 10 points
%
% =========================================================

\end{document}